\section{产业链关系矩阵}
光通信 A 股映射不能只分“模块和器件”。完整产业链至少包含八层：材料/基底、制造设备、光芯片/光源、光器件/无源、光模块、光纤光缆/ODN、网络设备、下游应用。本轮最强的短期收入弹性在 AI 光模块，最强的战略稀缺性在光芯片和关键器件，最稳定的周期底盘在光纤光缆与网络设备，最容易被忽视但决定良率的是制造设备和测试环节。
\begin{exhibitbox}[表：完整光通信产业链地图]
\centering\tiny
\renewcommand{\arraystretch}{1.18}
\setlength{\tabcolsep}{2pt}
\begin{tabularx}{\exhibitboxwidth}{>{\bfseries\raggedright\arraybackslash}p{1.25cm} >{\raggedright\arraybackslash\sloppy\hspace{0pt}}p{2.25cm} >{\raggedright\arraybackslash\sloppy\hspace{0pt}}p{2.25cm} >{\raggedright\arraybackslash\sloppy\hspace{0pt}}p{2.45cm} >{\raggedright\arraybackslash\sloppy\hspace{0pt}}X}
\toprule
\textbf{层级} & \textbf{核心对象} & \textbf{代表公司/映射} & \textbf{价值池} & \textbf{估值处理} \\
\midrule
材料/基底 & 高纯石英、光纤预制棒、InP/GaAs、薄膜铌酸锂、硅光晶圆、涂覆树脂 & 长飞光纤(601869)、亨通光电(600487)、中天科技(600522)；全球 Corning/材料厂 & 决定光纤衰减、芯片良率、调制器性能和长期供应安全 & A 股多数材料环节缺少纯标的，按产业瓶颈披露，不强行给目标价 \\
制造设备 & MOCVD/外延、光刻/刻蚀、划片、贴片、主动耦合、老化和光电测试 & 罗博特科观察池；海外设备/测试厂 & 决定高速光芯片、硅光、薄膜铌酸锂和模块封装良率 & 设备链重要但利润分母不稳定，作为观察池而非 PE 估值覆盖 \\
光芯片/光源 & DFB/EML/VCSEL/CW 激光器、PD/APD、硅光 PIC、PLC/AWG & 源杰科技(688498)、长光华芯(688048)、仕佳光子(688313) & 1.6T、硅光和 CPO 的战略稀缺层，国产替代弹性最高 & 给更高 PE 档位，但必须用客户认证和良率验证 \\
光器件/无源 & 隔离器、环形器、透镜、连接器、WDM、滤波器、精密耦合件 & 天孚通信(300394)、光库科技(300620)、太辰光(300570)、腾景科技(688195) & 影响模块良率、可靠性和 CPO/硅光封装难度 & 毛利率质量高，但估值依赖 attach rate 持续提升 \\
光模块 & 400G/800G/1.6T 可插拔模块、相干模块、光引擎 & 中际旭创(300308)、新易盛(300502)、长芯博创(300548)、德科立(688205)、联特科技(301205)、光迅科技(002281)、华工科技(000988)、剑桥科技(603083) & 云厂商 AI capex 到利润的最直接传导层 & 核心配置层，但价格已经提前反映高景气 \\
光纤/预制棒/海缆 & 预制棒、光纤、光缆、海缆和长距传输底座 & 长飞光纤(601869)、亨通光电(600487)、中天科技(600522) & 连接数据中心、运营商骨干、海缆和长距 DCI & 慢周期资产层，按现金流、项目交付和光纤价格定价 \\
通信线缆/ODN/网络集成 & 通信光缆、配线、ODN、通信线缆、网络工程和集成 & 通鼎互联(002491)、永鼎股份(600105)、特发信息(000070观察)、中贝通信(603220观察) & 承接 FTTH、运营商接入、园区网络和边缘节点 & 不套 AI 模块倍数；通鼎/永鼎按线缆与集成慢周期模型估值 \\
高速线缆/连接器/互连 & 高速数据线缆、射频同轴、连接器、铜互连和通信互连组件 & 兆龙互连(300913)、意华股份(002897)、神宇股份(300563)、欣天科技(300615观察) & 补足服务器、交换设备、园区和通信设备内部互连 & 按客户认证、毛利率和 Q2/Q3 利润兑现定价 \\
网络设备 & 交换机、路由器、OTN/ROADM、PON/接入设备、交换 ASIC/NIC 生态 & 中兴通讯(000063)、烽火通信(600498)、锐捷网络(301165观察)、光迅科技(002281)；海外 Broadcom/Marvell & 决定模块端口速率、相干传输和运营商投资节奏 & 业务更宽，需用折价倍数避免当作纯光模块估值 \\
下游应用 & AI 数据中心、DCI、运营商骨干、FTTH、企业园区、工业/汽车光互连 & 中兴通讯(000063)、烽火通信(600498)、共进股份(603118观察)、日海智能(002313观察)、天邑股份(300504观察)、模块/光纤全链条 & 决定需求周期长度：AI 短周期快，运营商与 FTTH 慢周期稳 & 报告将短周期弹性和慢周期底盘分开估值 \\
\bottomrule
\end{tabularx}
\end{exhibitbox}

\section{A 股全景股票图谱：覆盖池与观察池分开}
本报告将股票图谱分成两层：第一层是进入最终估值表的可估值覆盖池，必须具备当前价、正 EPS、股本、市值、三年 EPS 桥、目标价、空间和动作；第二层是产业链观察池，属于光通信或相邻网络链条，但因亏损、业务纯度不足、利润分母过薄或更偏服务/终端设备，暂不发布正式目标价。
\begin{exhibitbox}[表：可估值覆盖池 1]
\centering\tiny
\renewcommand{\arraystretch}{1.13}
\setlength{\tabcolsep}{2pt}
\begin{tabularx}{\exhibitboxwidth}{>{\bfseries\raggedright\arraybackslash}p{1.35cm} >{\raggedright\arraybackslash\sloppy\hspace{0pt}}p{2.3cm} >{\raggedright\arraybackslash\sloppy\hspace{0pt}}p{2.2cm} >{\centering\arraybackslash}p{0.75cm} >{\raggedright\arraybackslash\sloppy\hspace{0pt}}X}
\toprule
\textbf{标的} & \textbf{产业位置} & \textbf{层级} & \textbf{权重} & \textbf{纳入原因} \\
\midrule
中际旭创 300308 & AI 数据中心高速光模块 & 核心光模块 & 13\% & 正 EPS 且能建立当前价目标价模型；动作 中性 \\
新易盛 300502 & 高速光模块 & 核心光模块 & 12\% & 正 EPS 且能建立当前价目标价模型；动作 中性 \\
天孚通信 300394 & 精密光器件与封装平台 & 精密光器件 & 6\% & 正 EPS 且能建立当前价目标价模型；动作 中性观察 \\
光迅科技 002281 & 电信/数通光器件与光模块 & 光器件/光模块 & 4\% & 正 EPS 且能建立当前价目标价模型；动作 中性观察 \\
华工科技 000988 & 光模块与激光设备混合平台 & 混合光通信平台 & 4\% & 正 EPS 且能建立当前价目标价模型；动作 中性观察 \\
剑桥科技 603083 & 光模块与通信设备高 beta 标的 & 高弹性主题 & 2\% & 正 EPS 且能建立当前价目标价模型；动作 中性观察 \\
太辰光 300570 & 光纤连接器与无源光器件 & 无源器件 & 2\% & 正 EPS 且能建立当前价目标价模型；动作 中性观察 \\
源杰科技 688498 & 激光芯片与高速光源 & 战略上游 & 4\% & 正 EPS 且能建立当前价目标价模型；动作 中性观察 \\
亨通光电 600487 & 光纤光缆、海缆与电力通信一体化 & 光纤光缆 & 4\% & 正 EPS 且能建立当前价目标价模型；动作 中性观察 \\
中天科技 600522 & 光纤光缆、海缆与通信线缆 & 光纤光缆 & 4\% & 正 EPS 且能建立当前价目标价模型；动作 中性观察 \\
长飞光纤 601869 & 光纤预制棒、光纤与光缆 & 光纤/预制棒 & 4\% & 正 EPS 且能建立当前价目标价模型；动作 中性观察 \\
烽火通信 600498 & 电信光网络设备与系统 & 网络设备 & 3\% & 正 EPS 且能建立当前价目标价模型；动作 中性观察 \\
中兴通讯 000063 & 网络设备、运营商/云基础设施与终端应用 & 网络设备/应用 & 6\% & 正 EPS 且能建立当前价目标价模型；动作 中性 \\
\bottomrule
\end{tabularx}
\end{exhibitbox}
\begin{exhibitbox}[表：可估值覆盖池 2]
\centering\tiny
\renewcommand{\arraystretch}{1.13}
\setlength{\tabcolsep}{2pt}
\begin{tabularx}{\exhibitboxwidth}{>{\bfseries\raggedright\arraybackslash}p{1.35cm} >{\raggedright\arraybackslash\sloppy\hspace{0pt}}p{2.3cm} >{\raggedright\arraybackslash\sloppy\hspace{0pt}}p{2.2cm} >{\centering\arraybackslash}p{0.75cm} >{\raggedright\arraybackslash\sloppy\hspace{0pt}}X}
\toprule
\textbf{标的} & \textbf{产业位置} & \textbf{层级} & \textbf{权重} & \textbf{纳入原因} \\
\midrule
仕佳光子 688313 & PLC 分路器、AWG 与无源光芯片/器件 & 无源光芯片 & 3\% & 正 EPS 且能建立当前价目标价模型；动作 中性观察 \\
长光华芯 688048 & 高功率与数通激光芯片 & 激光芯片 & 2\% & 正 EPS 且能建立当前价目标价模型；动作 中性观察 \\
光库科技 300620 & 光纤器件、隔离器与铌酸锂/调制器期权 & 光器件 & 3\% & 正 EPS 且能建立当前价目标价模型；动作 减持 \\
通鼎互联 002491 & 光纤光缆、通信线缆与网络集成 & 光纤光缆 & 3\% & 正 EPS 且能建立当前价目标价模型；动作 中性观察 \\
永鼎股份 600105 & 光缆、通信线缆与电信基础设施集成 & 光纤光缆 & 3\% & 正 EPS 且能建立当前价目标价模型；动作 中性观察 \\
长芯博创 300548 & 光模块、有源光器件与相干传输组件 & 光模块/光器件 & 4\% & 正 EPS 且能建立当前价目标价模型；动作 中性观察 \\
德科立 688205 & 相干光模块与传输设备组件 & 相干/传输 & 3\% & 正 EPS 且能建立当前价目标价模型；动作 减持 \\
腾景科技 688195 & 光通信和激光用精密光学元件 & 精密光学 & 3\% & 正 EPS 且能建立当前价目标价模型；动作 减持 \\
联特科技 301205 & 高速光模块与数通收发器 & 高弹性光模块 & 2\% & 正 EPS 且能建立当前价目标价模型；动作 减持 \\
兆龙互连 300913 & 高速数据线缆、通信线缆与铜互连 & 高速线缆/互连 & 2\% & 正 EPS 且能建立当前价目标价模型；动作 减持 \\
意华股份 002897 & 高速连接器与通信互连组件 & 连接器/互连 & 2\% & 正 EPS 且能建立当前价目标价模型；动作 减持 \\
神宇股份 300563 & 射频同轴线缆与高速通信互连 & 线缆/互连 & 2\% & 正 EPS 且能建立当前价目标价模型；动作 减持 \\
\bottomrule
\end{tabularx}
\end{exhibitbox}
\begin{exhibitbox}[表：产业链观察池 1]
\centering\tiny
\renewcommand{\arraystretch}{1.13}
\setlength{\tabcolsep}{2pt}
\begin{tabularx}{\exhibitboxwidth}{>{\bfseries\raggedright\arraybackslash}p{1.35cm} >{\raggedright\arraybackslash}p{2.1cm} >{\raggedright\arraybackslash\sloppy\hspace{0pt}}p{3.0cm} >{\raggedright\arraybackslash\sloppy\hspace{0pt}}X}
\toprule
\textbf{代码} & \textbf{名称} & \textbf{产业位置} & \textbf{未纳入目标价原因} \\
\midrule
300757 & 罗博特科 & 光模块和硅光主动耦合/自动化设备 & 2026Q1 亏损；作为设备链观察池，不纳入当前价目标价覆盖。 \\
000070 & 特发信息 & 光纤光缆、通信设备与网络基础设施 & 当前数据包显示 2026Q1 亏损；利润分母修复前只进入观察池。 \\
300615 & 欣天科技 & 射频器件和通信连接件 & 光通信相关性存在，但产业链纯度和利润分母不足以支撑完整目标价。 \\
300565 & 科信技术 & 通信网络能源与机柜设备 & 更偏数据中心/通信配套，需另行验证订单和盈利质量。 \\
603118 & 共进股份 & 通信终端和网络设备代工 & 业务映射网络应用层，但光通信纯度和目标价分母不足。 \\
002313 & 日海智能 & 通信网络服务和物联网平台 & 更偏系统集成与服务，不适合直接套光通信硬件估值。 \\
\bottomrule
\end{tabularx}
\end{exhibitbox}
\begin{exhibitbox}[表：产业链观察池 2]
\centering\tiny
\renewcommand{\arraystretch}{1.13}
\setlength{\tabcolsep}{2pt}
\begin{tabularx}{\exhibitboxwidth}{>{\bfseries\raggedright\arraybackslash}p{1.35cm} >{\raggedright\arraybackslash}p{2.1cm} >{\raggedright\arraybackslash\sloppy\hspace{0pt}}p{3.0cm} >{\raggedright\arraybackslash\sloppy\hspace{0pt}}X}
\toprule
\textbf{代码} & \textbf{名称} & \textbf{产业位置} & \textbf{未纳入目标价原因} \\
\midrule
603220 & 中贝通信 & 通信网络建设和算力基础设施服务 & 项目制属性强，需以订单、现金流和负债结构另行建模。 \\
300504 & 天邑股份 & 通信终端、网关和接入设备 & 接入设备相关，但 AI 光通信弹性有限。 \\
002902 & 铭普光磁 & 通信磁性器件和光电模块配套 & 配套属性存在，但需要更清晰收入占比和利润分母。 \\
301165 & 锐捷网络 & 企业网络设备和交换路由产品 & 网络设备应用层标的，可后续纳入独立网络设备报告。 \\
GLOBAL & MOCVD / lithography / etch / test equipment & 石英、InP/GaAs、铌酸锂、硅光晶圆等材料底座 & 材料方向纳入链条图谱；可投资标的需要单独材料公司报告。 \\
GLOBAL & Quartz / InP / GaAs / LiNbO3 / silicon photonics wafers & 石英、InP/GaAs、铌酸锂、硅光晶圆等材料底座 & 材料方向纳入链条图谱；可投资标的需要单独材料公司报告。 \\
\bottomrule
\end{tabularx}
\end{exhibitbox}

\section{产业链层级深挖：每一层该问什么}
\subsection{材料/基底：决定物理边界，不直接等于短期弹性}
材料层是光通信的底层约束。光纤预制棒和高纯石英决定长距离传输的衰减、成本和供给弹性；InP/GaAs 决定高速激光器和探测器的外延质量；薄膜铌酸锂和硅光晶圆决定调制器、PIC 和 CPO 路线的长期选择。材料环节的商业特征是技术壁垒高、验证周期长、收入确认慢，短期股价往往跟随主题扩散，但利润兑现并不一定同步。

调研上不能只问“有没有供货”，而要看三类指标：第一，材料纯度、缺陷密度、尺寸和批间稳定性是否满足高速光通信要求；第二，客户认证是否进入量产而不是样品；第三，价格和良率改善是否足以抵消扩产折旧。A 股能直接建模的纯材料标的有限，因此本报告把材料作为产业链必备层披露，但不把缺少稳定利润分母的对象放入最终目标价表。
\subsection{制造设备：良率瓶颈会先于利润表出现}
高速光芯片、硅光、薄膜铌酸锂和 1.6T 模块的放量，都离不开外延、光刻/刻蚀、划片、贴片、主动耦合、老化和光电测试设备。设备不是“配角”，而是良率和吞吐的前置条件。尤其在 1.6T 和 CPO 过渡期，主动耦合精度、热稳定性、自动化节拍和高速测试一致性会决定单线产能和返工率。

但设备公司进入投资覆盖需要更严格：设备订单可能领先收入几个季度，客户认证也可能停在小批量，单季亏损或项目制收入会让 PE 模型失真。罗博特科这类标的适合进入观察池，重点跟踪新签订单、客户结构、交付节奏、毛利率、费用率和经营现金流，而不是在亏损分母上硬推目标价。
\subsection{光芯片/光源：战略稀缺性最高，验证门槛也最高}
光芯片层包括 DFB/EML/VCSEL/CW 激光器、PD/APD、PLC/AWG、硅光 PIC 和调制器，是 1.6T、硅光和 CPO 的核心期权。源杰科技、长光华芯、仕佳光子等标的的价值不在于短期收入规模，而在于能否进入头部模块厂和系统厂供应链，并在良率、功耗、可靠性和成本上通过量产考验。

估值上，光芯片可以给战略溢价，但不能忽略利润分母。若客户认证停留在送样阶段，收入占比低或毛利率不稳定，高 PE 会迅速变成估值陷阱。本报告对上游芯片设置更高倍数区间，同时也把证据质量、Q2/Q3 放量和客户数列为核心失效条件。
\subsection{光器件/无源：毛利率质量高，但需要 attach rate 支撑}
光器件层包括隔离器、环形器、透镜、WDM、滤波器、连接器、精密耦合件和部分调制器件。天孚通信、光库科技、太辰光的共同点是毛利率质量较好，且受益于高速模块封装复杂度上升。不同点在于客户结构、产品稀缺性和订单兑现节奏。

调研重点不是“有没有 AI 逻辑”，而是 attach rate 和价值量是否真的上升。如果高速模块单机价值量提升被价格下降抵消，或客户开始自供关键器件，器件公司的利润弹性会低于主题预期。因此本报告要求用单季毛利率、收入增速和经营现金流三项共同验证。
\subsection{光模块：最直接的 AI capex 弹性，也是最拥挤的交易}
光模块是 AI 数据中心向 A 股传导最直接的利润层。中际旭创和新易盛在 2026Q1 已经把海外 AI 订单转化为收入、利润和较高毛利率，因此它们仍是核心研究池。光迅科技、华工科技、剑桥科技也有模块业务，但业务混合度、利润率和客户结构决定了它们不能完全按纯 AI 模块龙头估值。

模块层的关键变量包括 800G 出货持续性、1.6T 资格认证、ASP 下行幅度、客户集中、汇率、海外产能和现金流。若收入增长但应收账款、存货和经营现金流恶化，说明订单质量弱于利润表表象。目标价必须随 EPS 和毛利率动态调整，不能沿用主题行情高点的静态倍数。
\subsection{光纤光缆/ODN：慢周期底盘，不是 AI 模块替代品}
光纤光缆和 ODN 连接数据中心、城域 DCI、运营商骨干、FTTH、海缆和边缘节点。长飞光纤、亨通光电、中天科技承担预制棒、光纤光缆和海缆底座；通鼎互联、永鼎股份更偏通信线缆、ODN、网络集成和接入侧场景。它们的利润弹性通常慢于 AI 模块，且受价格周期、运营商招标、海缆项目、FTTH/园区项目和库存周期影响更大。

投资处理上，这一层适合用现金流、订单能见度和项目利润率估值，而不是套用高端光模块 PE。若光纤价格改善、预制棒利用率恢复、海缆/运营商项目交付提速，慢周期底盘可以提供组合稳定性；若价格竞争重新加剧，它们会拖累全链条加权空间。
\subsection{网络设备：把光通信从器件拉回系统和应用}
网络设备层包括交换机、路由器、OTN、ROADM、PON、接入设备以及与 switch ASIC/NIC 配套的系统生态。中兴通讯和烽火通信把光通信连接到运营商、云网络和企业应用场景。它们的优势是业务规模和客户稳定性，劣势是光通信纯度低、周期长、项目利润率波动。

这一层的调研要看运营商 capex、招标份额、网络设备毛利率、云网络产品进展和海外政策压力。估值上，网络设备公司应给更低但更稳的倍数，除非其 AI 网络产品能够证明新的利润曲线。
\subsection{下游应用：决定估值应该快还是慢}
下游应用不是一个变量。AI 数据中心是快周期，几个季度内就会体现在模块订单和毛利率；DCI/相干传输介于云和运营商之间；运营商骨干、FTTH 和企业园区更偏慢周期；工业和汽车光互连商业化更分散。把这些需求混成一个“光通信景气”会导致估值错配。

因此本报告把 25 只覆盖标的按应用周期分层：模块龙头看 AI 快周期，光纤光缆看运营商/海缆/FTTH，网络设备看 capex 和利润率，上游芯片/器件看认证与良率。只有当多个应用同时兑现，产业链整体才适合上修；若只有 AI 模块兑现，其他层级不应跟随无差别追高。

\section{设备与材料为什么单列观察池}
设备和材料是产业链必需项，但并不等于每个环节都适合纳入本报告最终估值表。估值表要求当前价、股本、市值、2026E/2027E/2028E EPS、目标价、合理区间和空间均可计算；如果公司当期亏损、业务混杂、或 A 股缺少纯标的，本报告只做观察池披露，避免为了“覆盖完整”而制造虚假的目标价。
\begin{exhibitbox}[表：设备与材料观察池]
\centering\scriptsize
\begin{tabularx}{\exhibitboxwidth}{>{\bfseries\raggedright\arraybackslash}p{1.4cm} >{\raggedright\arraybackslash}p{2.2cm} >{\raggedright\arraybackslash\sloppy\hspace{0pt}}X >{\raggedright\arraybackslash\sloppy\hspace{0pt}}X}
\toprule
\textbf{代码} & \textbf{名称} & \textbf{产业位置} & \textbf{处理方式} \\
\midrule
300757 & 罗博特科 & 光模块和硅光主动耦合/自动化设备 & 2026Q1 亏损；作为设备链观察池，不纳入当前价目标价覆盖。 \\
000070 & 特发信息 & 光纤光缆、通信设备与网络基础设施 & 当前数据包显示 2026Q1 亏损；利润分母修复前只进入观察池。 \\
300615 & 欣天科技 & 射频器件和通信连接件 & 光通信相关性存在，但产业链纯度和利润分母不足以支撑完整目标价。 \\
300565 & 科信技术 & 通信网络能源与机柜设备 & 更偏数据中心/通信配套，需另行验证订单和盈利质量。 \\
603118 & 共进股份 & 通信终端和网络设备代工 & 业务映射网络应用层，但光通信纯度和目标价分母不足。 \\
002313 & 日海智能 & 通信网络服务和物联网平台 & 更偏系统集成与服务，不适合直接套光通信硬件估值。 \\
603220 & 中贝通信 & 通信网络建设和算力基础设施服务 & 项目制属性强，需以订单、现金流和负债结构另行建模。 \\
300504 & 天邑股份 & 通信终端、网关和接入设备 & 接入设备相关，但 AI 光通信弹性有限。 \\
002902 & 铭普光磁 & 通信磁性器件和光电模块配套 & 配套属性存在，但需要更清晰收入占比和利润分母。 \\
301165 & 锐捷网络 & 企业网络设备和交换路由产品 & 网络设备应用层标的，可后续纳入独立网络设备报告。 \\
GLOBAL & MOCVD / lithography / etch / test equipment & 石英、InP/GaAs、铌酸锂、硅光晶圆等材料底座 & 材料方向纳入链条图谱；可投资标的需要单独材料公司报告。 \\
GLOBAL & Quartz / InP / GaAs / LiNbO3 / silicon photonics wafers & 石英、InP/GaAs、铌酸锂、硅光晶圆等材料底座 & 材料方向纳入链条图谱；可投资标的需要单独材料公司报告。 \\
\bottomrule
\end{tabularx}
\end{exhibitbox}

\section{竞争格局}
中际旭创和新易盛的优势来自规模、海外客户认证和高端模块交付经验；天孚通信、光库科技、太辰光的优势来自精密组件、封装和无源器件；源杰科技、长光华芯、仕佳光子的优势在光芯片、激光器和 PLC/AWG 等上游稀缺环节；长飞光纤、亨通光电、中天科技构成光纤光缆和海缆底盘；通鼎互联、永鼎股份补足通信线缆、ODN 和网络集成慢周期敞口；兆龙互连、意华股份、神宇股份补足高速线缆和连接器；中兴通讯、烽火通信和光迅科技则把光通信接入网络设备和运营商/云网络场景。估值上，不能把上述层级统一套一个“AI 光模块倍数”：利润弹性、客户集中、资本开支周期和可替代性都不同。
